\section{Discussion}

\textbf{Outline for review, to be written as prose (target about 1,000-1,200 words).} Each point gives the claim, the evidence in this paper, how it relates to prior work, and what stays open.

\begin{itemize}
\item
  \textbf{D1. The judge decides, and precision matters more than volume.}
  \emph{Claim:} what decided a life was the share of failed attempts among its lessons, not their number or the kind of judge.
  \emph{Evidence:} $\rho$ = -0.81 across 97 lives (5.6); with sense v1 and the \emph{stated verdict}, seven of the eight lives that learned at least once reached 34-61 from only 1-7 lessons; one failed lesson on night 1 ended six of seven sense lives.
  \emph{Prior work:} the same lesson as self-training on verified successes \cite{zelikman2022star, gulcehre2023reinforced, singh2024beyond}, and the damage of training on wrong or self-generated data \cite{shumailov2024ai}.
  \emph{Implication:} a judge should abstain when in doubt; early lessons matter most, because there are few of them.
\item
  \textbf{D2. Why self-judging decays.}
  \emph{Claim:} a model that judges its own work learns at first and then declines, and neither time nor the adapter design changes this.
  \emph{Evidence:} Experiments 1-A, 1-B (a year) and 1-C (four adapters); the self-judged lives let in 17-47\% failures (5.6).
  \emph{Prior work:} models cannot reliably correct themselves without outside feedback \cite{huang2024large}; model judges are biased \cite{zheng2023judging}; optimising a proxy reward eventually lowers the true one \cite{gao2023scaling}. Self-rewarding models report gains over a few iterations \cite{yuan2024selfrewarding, wu2025meta}; our longer horizon and much smaller models are a different setting, so this is a caution, not a contradiction.
\item
  \textbf{D3. Discernment: telling politeness from approval.}
  \emph{Claim:} a separate, frozen judge that reads the user avoids the learner grading itself, but it must tell a clear verdict on the result from thanks for the effort.
  \emph{Evidence:} v1 worked with a plain user (5.2); v2 let polite give-ups in and 3 of 5 lives collapsed (5.3); v3, which first asks whether the reaction is a clear verdict, cut give-ups and polite complaints read as success from 15.8\% to 1.0\%, and four of its five lives stayed near the \emph{oracle} for two weeks (5.4). The teacher itself read half of the politely unhappy and sarcastic reactions as a yes (3.2).
  \emph{Prior work:} implicit feedback can be more useful than explicit ratings \cite{kreutzer2018can}, is common in real chats \cite{doneyehiya2024naturally, shi2026wildfeedback} and can improve models over rounds \cite{hancock2019learning, chen2025retrospective}; we add its main failure mode when learning runs unattended: polite reactions to failure.
  \emph{Cost:} v3 took fewer real approvals (59\% against 70\%), so its lessons were scarce (D5).
\item
  \textbf{D4. Discernment is bounded by what the user can see.}
  \emph{Claim:} no reading of a reaction can recover what the user did not know.
  \emph{Evidence:} on $\tau^2$-bench the customer cannot see the database and sincerely thanked failed tasks; sense v3 read them correctly, so 80\% of its lessons were failures and the 1.7B lives collapsed (5.5).
  \emph{Prior work:} in the work we cite on simulated users \cite{yao2025taubench, barres2026tau2, ai2026memorybench} and on learning after deployment \cite{hancock2019learning, xu2023learning}, what the user can verify is not treated separately from what they say.
  \emph{Implication:} trust a verdict only where the user can check the result; elsewhere, mix in verifiable signals.
\item
  \textbf{D5. Correct lessons are not enough: scarcity and imbalance.}
  \emph{Claim:} clean lessons can still fade, and correct lessons can still cause forgetting.
  \emph{Evidence:} sense v3's four surviving lives, at 2-10\% failed lessons and about 2 lessons a day, fell from near the \emph{oracle} on day 14 to 9-22 by day 60 (5.4); on $\tau^2$-bench, three quarters of the \emph{oracle}'s lessons came from airline, and the 1.7B \emph{oracle} lives rose on airline while retail and telecom fell (5.5).
  \emph{Prior work:} forgetting under continual fine-tuning and replay as its standard remedy \cite{mccloskey1989catastrophic, rolnick2019experience, delange2022continual, luo2025empirical}; LoRA limits forgetting but does not remove it \cite{biderman2024lora}; consolidation in a separate sleep phase \cite{mcclelland1995why, behrouz2026language} is another route.
  \emph{Open:} why the fade happens; whether replay balanced by domain or task type prevents the forgetting.
\item
  \textbf{D6. What to do, and what comes next.} A short closing paragraph.
  \emph{Guidance:} trust a verdict only when the user can check the result; abstain on unclear reactions rather than guess; mix in verified outcomes where available; balance replay across skills; for a model on a device, keep the judge small too, as the sense is \cite{liu2024mobilellm, zhu2023pockengine}.
  \emph{Next:} an \emph{oracle} limited to the sense's lesson rate, to explain the fade; balanced replay on $\tau^2$-bench; the sense on real users' reactions; larger models (Section 7); and a fair comparison of the nightly loop with offline training and RL on the same verified reward, on a standard benchmark split \cite{song2024trial, singh2024beyond}.
\end{itemize}
